\section{Independent ownership of target multigrid and source DF}
\label{sec:target-multigrid}

\subsection{Two independent choices}

A periodic DC-DFT calculation contains two numerical choices with different
roles.  The \emph{source two-electron backend} evaluates Coulomb, exchange,
and orbital-Hessian actions for the source HF problem.  This backend is
FFTDF, GDF, or RSDF.  The \emph{target integrator} evaluates target $J$ and
$E_{\mathrm{xc}}$ at the converged HF density.  Calling both choices one
``DF method'' obscures their separate responsibilities.

We denote a route by the ordered pair
\begin{equation}
  \mathcal R
  =\left(\mathcal B_{\mathrm{source}},\mathcal I_{\mathrm{target}}\right),
  \qquad
  \mathcal B_{\mathrm{source}}\in\{\mathrm{FFTDF},\mathrm{GDF},\mathrm{RSDF}\}.
  \label{eq:target-owner-split}
\end{equation}
Changing the target integrator must not change the source HF density
$D^{\mathrm{HF}}$, HF energy, occupied and virtual orbitals, or source Hessian
$A^{\mathrm{HF}}$.  It changes the target quadrature in
$E_{\mathrm{DFA}}[D^{\mathrm{HF}}]$, the target Fock matrix, the orbital-gradient
right-hand side, and the explicit target derivative.

\subsection{Three target-integration modes}

The current GPU route distinguishes three modes.

\begin{center}
\begin{longtable}{p{0.21\textwidth}p{0.22\textwidth}p{0.22\textwidth}p{0.23\textwidth}}
\toprule
Mode & Target $J$ & Target XC & Source HF/AX \\
\midrule
\texttt{uniform} & selected DF & uniform KNumInt & selected DF \\
\texttt{multigrid\_xc} & selected DF & multigrid & selected DF \\
\texttt{multigrid\_j\_xc} & multigrid & multigrid & selected DF \\
\bottomrule
\end{longtable}
\end{center}

For \texttt{multigrid\_xc}, called M1, the target effective potential is
\begin{equation}
  V_{\mathrm{eff}}^{\mathrm{M1}}[D]
  =J_{\mathcal B_{\mathrm{source}}}[D]
   +V_{\mathrm{xc}}^{\mathrm{MG}}[D].
  \label{eq:target-m1-veff}
\end{equation}
M1 keeps the finite-auxiliary GDF/RSDF representation of target $J$ and
replaces only the XC quadrature.

For \texttt{multigrid\_j\_xc}, called M2,
\begin{equation}
  V_{\mathrm{eff}}^{\mathrm{M2}}[D]
  =J_{\mathrm{MG}}[D]+V_{\mathrm{xc}}^{\mathrm{MG}}[D].
  \label{eq:target-m2-veff}
\end{equation}
The multigrid $J_{\mathrm{MG}}$ and a finite-auxiliary GDF/RSDF $J$ are
different finite numerical representations.  They approach the same Coulomb
limit when their meshes and auxiliary bases are separately converged, but
their finite-setting energies need not agree.  M2 is therefore a separate
target model, not merely a drop-in acceleration of GDF or RSDF.

\subsection{Ownership in the energy and gradient}

The target Coulomb derivative must be included exactly once.  The ownership
rule is
\begin{equation}
  \mathbf g_{\mathrm{target}}
  =\begin{cases}
    \mathbf g_J^{\mathcal B_{\mathrm{source}}}
    +\mathbf g_{\mathrm{xc}}^{\mathrm{uniform}},
      &\mathcal I=\texttt{uniform},\\
    \mathbf g_J^{\mathcal B_{\mathrm{source}}}
    +\mathbf g_{\mathrm{xc}}^{\mathrm{MG}},
      &\mathcal I=\texttt{multigrid\_xc},\\
    \mathbf g_{J+\mathrm{xc}}^{\mathrm{MG}},
      &\mathcal I=\texttt{multigrid\_j\_xc}.
  \end{cases}
  \label{eq:target-gradient-owner}
\end{equation}
Adding the selected-DF target $dJ/dR$ in the last line double counts the
Coulomb response.  Omitting it in M1 differentiates a Hamiltonian different
from the one used for the energy.

The target hcore and GTH pseudopotential remain owned by the selected
one-electron provider in all three modes.  Attaching a multigrid object to the
entire target KS object can implicitly redirect upstream \texttt{get\_hcore()}
or \texttt{get\_j()} dispatch.  The implementation must instead combine the
low-level target XC/J outputs explicitly and report separate hcore, PP, J,
and XC owners.

\subsection{Why target multigrid does not replace the HF Hessian}

The DC-DFT adjoint equation differentiates the source HF stationarity
constraint.  In
\begin{equation}
  A_{\mathrm{HF}}^Tz=-g_{\mathrm{target}},
\end{equation}
the right-hand side depends on the target integrator, whereas
$A_{\mathrm{HF}}$ remains the source Hessian.  Replacing AX J/K by multigrid
after selecting M1 or M2 would solve the response of a different stationarity
constraint than the one satisfied by the source density.

The workflow is therefore:
\begin{enumerate}
  \item converge source HF with GDF or RSDF and obtain $D^{\mathrm{HF}}$;
  \item build $E_{\mathrm{DFA}}[D^{\mathrm{HF}}]$ and $F_{\mathrm{target}}$
        with the selected target integrator;
  \item form $g_{\mathrm{target}}$ from $F_{\mathrm{target},vo}$;
  \item solve for $z$ with source GDF/RSDF AX;
  \item assemble target explicit and source-HF response derivatives.
\end{enumerate}

\subsection{Source and target meshes}

The \texttt{cell.mesh} resolved from \texttt{cell.ke\_cutoff} may be used by
FFT-based source operators, the GTH local potential, and other reciprocal
operations.  The target multigrid mesh is an independent parameter for target
XC/J quadrature.  Equal default values do not imply equal accuracy meaning.

Convergence should be separated as
\begin{equation}
  \lim_{\mathcal M_s\to\infty}
  \lim_{\mathcal M_t\to\infty}
  \mathbf g(\mathcal M_s,\mathcal M_t).
  \label{eq:target-grid-convergence}
\end{equation}
Reporting only one cutoff does not identify which representation is
converged.  \texttt{target\_\allowbreak multigrid\_\allowbreak mesh=null}
inherits the current cell
mesh; it is not an automatic memory cap.  Large-cell reports must record the
resolved triplet and number of grid points.

Multigrid distributes work over Gaussian/locality levels, but it remains a
finite quadrature.  It is not automatically more accurate or more
translation invariant than a dense uniform grid.  Energy phase scans,
translation sums, and complete-gradient mesh convergence test distinct
properties and should all be reported.

\subsection{Workspace reuse and lifetime}

Multigrid tasks, sorted Gaussian pairs, and the target Fock built during the
target energy can be reused by the gradient.  Using a different multigrid
object or mesh in the gradient does not differentiate the central energy.
Conversely, keeping the large task cache after its final use overlaps its
memory with later GDF/RSDF response work.

The safe lifetime is
\begin{equation}
  \text{build at target energy}
  \longrightarrow
  \text{reuse at target gradient}
  \longrightarrow
  \text{release before DF response}.
  \label{eq:target-multigrid-memory}
\end{equation}
This is ownership management after the last consumer, not an approximation
that discards grid points.

\subsection{Recommended scope}

M1 is the conservative opt-in route because it retains source HF and the
finite-auxiliary target $J$ while replacing only target XC.  M2 can be
analytically consistent, but it also changes the target-$J$ representation
and requires separate convergence.  A PP-only multigrid route is inconsistent
until central hcore, explicit PP derivatives, the nuclear CPHF source, and the
Z contraction all use the same representation.
